\section{第3组：刚体定轴转动}
\sourceproblem{3-T1}{9}{13}
\problemtext 轮从静止经6秒匀加速至200转每分钟，匀速一段时间后制动，5分钟后停止，总转数3100，求总时间。
\workings 若制动为匀减速，则加速段 $N_1=100(6/60)=10$ 转，制动段 $N_3=100(5)=500$ 转，匀速段 $N_2=2590$ 转、耗时 $2590/200=12.95$ 分钟。总时间 $18.05$ 分钟，即1083秒。
\crosscheck 原题只明确加速段均匀增加，未明确制动规律；若制动非匀减速，$N_3$ 无法由端点转速和历时唯一确定，总时间也不唯一。所给数值以补充匀减速条件为前提。
\sourceproblem{3-T2}{9}{13}
\problemtext 半径1.3米圆周运动，初始 $\theta=\omega=0$，角加速度 $\beta=120t^2-48t+16$（SI），求角速度、角位置及切法向加速度。
\workings 积分得
\[
\omega=40t^3-24t^2+16t,\qquad \theta=10t^4-8t^3+8t^2.
\]
切向分量 $a_\tau=1.3(120t^2-48t+16)$，法向分量 $a_n=1.3(40t^3-24t^2+16t)^2$（SI）。
\crosscheck 两次求导恢复角加速度与初值；法向分量非负，指向圆心。
\sourceproblem{3-T3}{9}{13}
\problemtext 重物 $m$ 由静止拉动惯量 $J=MR^2/2$ 的定滑轮，求下落速度随时间的关系。
\workings $mg-T=ma,\ TR=J a/R$，故 $a=mg/(m+M/2)$，下落速率 $v(t)=mgt/(m+M/2)$。
\crosscheck 能量式 $mgh=(m+M/2)v^2/2$ 与匀加速关系一致；无滑动且轴光滑。
\sourceproblem{3-T4}{10}{14}
\problemtext 飞轮轴半径2.00厘米，20.0克物体匀速下落；改挂5.00千克物体从静止10.0秒下落40.0厘米，取 $g=9.80$，求轴摩擦力矩、惯量及第二次张力。
\workings 若摩擦力矩大小恒定，首次 $\tau_f=mg r=0.00392\,\mathrm{N\,m}$。第二次 $a=2h/t^2=0.00800\,\mathrm{m/s^2}$，$T=M(g-a)=48.96\,\mathrm N$，由 $Tr-\tau_f=J a/r$ 得
\[
J=(Tr-\tau_f)r/a=2.4382\,\mathrm{kg\,m^2}.
\]
按有效数字 $J\approx2.44\,\mathrm{kg\,m^2}$，$T\approx49.0\,\mathrm N$。
\crosscheck 两次实验须具有相同的恒定轴摩擦力矩，且绳不打滑；若摩擦随载荷或转速变化，题设不足以确定惯量。
\sourceproblem{3-T5}{10}{14}
\problemtext 两等质量 $m$、等半径 $r$ 均质滑轮与两端重物 $m,2m$ 组成系统，求两轮间张力。
\workings 取 $2m$ 端向下，整个系统等效惯性质量为 $m+2m+2(J/r^2)=4m$，故 $a=g/4$。轻端张力 $T_1=m(g+a)$，中间张力满足 $(T-T_1)r=J a/r$，得
\[
T=m(g+a)+ma/2=11mg/8.
\]
重端张力 $T_2=2m(g-a)=3mg/2$。
\crosscheck 另一滑轮满足 $T_2-T=ma/2$，与左轮一致；三段张力不可视为相等。
\sourceproblem{3-T6}{11}{15}
\problemtext 转台初始每秒1转，总惯量6.0，收臂后2.0（$\mathrm{kg\,m^2}$），求末角速率及转动动能增量。
\workings 初角速率 $\omega_i=2\pi\,\mathrm{rad/s}$。角动量守恒给出 $\omega_f=3\omega_i=6\pi\,\mathrm{rad/s}$（每秒3转）。动能差
\[
\Delta K=\tfrac12(2)(6\pi)^2-\tfrac12(6)(2\pi)^2=24\pi^2\,\mathrm J\approx237\,\mathrm J.
\]
\crosscheck $K=L_z^2/(2J)$，惯量减至三分之一使转动动能成为原来的三倍；增量来自人体做功。
\sourceproblem{3-T7}{11}{15}
\problemtext 均质细杆质量 $m$、长 $l$，端点惯量 $ml^2/3$，水平静止释放，求与向下竖直方向成30度时角速率。
\workings 质心下降 $(l/2)\cos30^\circ$，故 $mg(l/2)\cos30^\circ=(ml^2/6)\omega^2$，得 $\omega=\sqrt{3\sqrt3g/(2l)}$。
\crosscheck 初始水平时角速率零，最低点速率对应 $\omega=\sqrt{3g/l}$；本题答案介于两者之间。
\sourceproblem{3-T8}{11}{15}
\problemtext 自由旋转圆盘惯量 $J$、半径 $R$，边缘有质量 $m$ 的人，初角速率 $\omega$；人走到盘心，求角速度和动能变化。
\workings 初末惯量为 $J+mR^2,J$，故 $\omega_f=(1+mR^2/J)\omega$，
\[
\Delta\omega=mR^2\omega/J,\qquad
\Delta K=\frac{mR^2(J+mR^2)}{2J}\omega^2.
\]
\crosscheck 取人为质点，初末均相对盘静止；轴向角动量守恒，转动动能增加，不主张机械能守恒。

\sourceproblem{3-T9}{12}{16}
\problemtext 均质杆长0.4米、质量1千克，绕上端轴静止下垂；8克小物体以200米每秒水平嵌入距轴 $3l/4$ 处，求碰后角速度、最大偏角。
\workings 设 $d=3l/4$，$J=Ml^2/3+md^2$。碰撞瞬间 $\omega=mvd/J$，约 $8.88\,\mathrm{rad/s}$。随后无耗散摆动，
\[
\frac12J\omega^2=g(Ml/2+md)(1-\cos\theta_{\max}).
\]
取 $g=9.8\,\mathrm{m/s^2}$，得 $\cos\theta_{\max}\approx-0.0746$，$\theta_{\max}\approx94.3^\circ$。
\crosscheck 碰撞阶段只守恒轴向角动量，摆动阶段才守恒机械能；转动能不足以到达倒立位置。
\sourceproblem{3-T10}{12}{16}
\problemtext 长 $l$、质量 $M$ 的均质杆从与竖直成 $\theta_0$ 的位置静止下摆，在最低点与光滑地面上的静止质点 $m$ 弹性碰撞，求碰后两者速度。
\workings 端轴惯量 $J=Ml^2/3$，碰前角速率 $\omega_0=\sqrt{3g(1-\cos\theta_0)/l}$。取杆端碰前运动方向为正，碰撞满足
\[
J\omega_0=J\omega+mlv,\qquad v-l\omega=l\omega_0.
\]
解得 $\omega=(M-3m)\omega_0/(M+3m)$，$v=2Ml\omega_0/(M+3m)$。
\crosscheck 恢复系数为1；代回能量守恒。$m>M/3$ 时杆反向，不能把有向角速度强行取正。
